\documentclass[10pt,a4paper]{amsart}
\usepackage{amsmath,amssymb,amsthm,graphicx}
\usepackage{verbatim}
\newenvironment{reference}{\comment}{\endcomment}
%\newenvironment{reference}{}{}
\newenvironment{slogan}{\comment}{\endcomment}
\newenvironment{history}{\comment}{\endcomment}

% For commutative diagrams we use Xy-pic
\usepackage[all]{xy}

% We use 2cell for 2-commutative diagrams.
\xyoption{2cell}
% Theorem environments.
%
\theoremstyle{plain}
\newtheorem{theorem}[subsection]{Theorem}
\newtheorem{proposition}[subsection]{Proposition}
\newtheorem{lemma}[subsection]{Lemma}

\theoremstyle{definition}
\newtheorem{definition}[subsection]{Definition}
\newtheorem{example}[subsection]{Example}
\newtheorem{exercise}[subsection]{Exercise}
\newtheorem{situation}[subsection]{Situation}

\theoremstyle{remark}
\newtheorem{remark}[subsection]{Remark}
\newtheorem{remarks}[subsection]{Remarks}

\numberwithin{equation}{subsection}

% Macros
%
\def\lim{\mathop{\mathrm{lim}}\nolimits}
\def\colim{\mathop{\mathrm{colim}}\nolimits}
\def\Spec{\mathop{\mathrm{Spec}}}
\def\Hom{\mathop{\mathrm{Hom}}\nolimits}
\def\Ext{\mathop{\mathrm{Ext}}\nolimits}
\def\SheafHom{\mathop{\mathcal{H}\!\mathit{om}}\nolimits}
\def\SheafExt{\mathop{\mathcal{E}\!\mathit{xt}}\nolimits}
\def\Sch{\mathit{Sch}}
\def\Mor{\mathop{\mathrm{Mor}}\nolimits}
\def\Ob{\mathop{\mathrm{Ob}}\nolimits}
\def\Sh{\mathop{\mathit{Sh}}\nolimits}
\def\NL{\mathop{N\!L}\nolimits}
\def\CH{\mathop{\mathrm{CH}}\nolimits}
\def\proetale{{pro\text{-}\acute{e}tale}}
\def\etale{{\acute{e}tale}}
\def\QCoh{\mathit{QCoh}}
\def\Ker{\mathop{\mathrm{Ker}}}
\def\Im{\mathop{\mathrm{Im}}}
\def\Coker{\mathop{\mathrm{Coker}}}
\def\Coim{\mathop{\mathrm{Coim}}}

% Boxtimes
%
\DeclareMathSymbol{\boxtimes}{\mathbin}{AMSa}{"02}

%
% Macros for moduli stacks/spaces
%
\def\QCohstack{\mathcal{QC}\!\mathit{oh}}
\def\Cohstack{\mathcal{C}\!\mathit{oh}}
\def\Spacesstack{\mathcal{S}\!\mathit{paces}}
\def\Quotfunctor{\mathrm{Quot}}
\def\Hilbfunctor{\mathrm{Hilb}}
\def\Curvesstack{\mathcal{C}\!\mathit{urves}}
\def\Polarizedstack{\mathcal{P}\!\mathit{olarized}}
\def\Complexesstack{\mathcal{C}\!\mathit{omplexes}}
% \Pic is the operator that assigns to X its picard group, usage \Pic(X)
% \Picardstack_{X/B} denotes the Picard stack of X over B
% \Picardfunctor_{X/B} denotes the Picard functor of X over B
\def\Pic{\mathop{\mathrm{Pic}}\nolimits}
\def\Picardstack{\mathcal{P}\!\mathit{ic}}
\def\Picardfunctor{\mathrm{Pic}}
\def\Deformationcategory{\mathcal{D}\!\mathit{ef}}
\usepackage{cleveref}
\crefname{originaldefinition}{original definition}{original definitions}
\Crefname{originaldefinition}{Original definition}{Original definitions}
\makeatletter
\newlabel{spaces-definition-quotient}{{02Z3}{source}{}{sourcepointer-1}{}}
\newlabel{spaces-definition-quotient@cref}{{[originaldefinition][0][]02Z3}{[1][0][]source}{}{}{}}
\makeatother
\title{An etale morphism of algebraic spaces whose fibres contain nontrivial specializations}
\author{The Stacks Project Authors}
\date{stacks-project:examples.tex:local-snapshot-20261012; GNU FDL 1.2 or later}
\begin{document}
\maketitle
\noindent\textit{Editorial note.} All delivered mathematical text is the author's own, copied from the pinned Stacks Project chapter file examples.tex; no statement or proof step was rewritten, reordered or removed.\par\medskip
\noindent\textit{Editorial note.} The counted proof is the author's complete argument: The specialization chain lives in |U| (the etale cover), which is exactly the shape the counted statement asks for: specializations inside a fibre of an etale morphism of algebraic spaces.\par\medskip
\noindent\textit{Editorial note.} The label lemma-specializations-fibre-etale names the counted statement, and the proof environment immediately following it in the source is delivered verbatim as its proof.\par\medskip
\noindent\textit{Editorial note.} The counted statement is carried by the author's own lemma environment in the chapter source.\par\medskip
\noindent\textit{Allowed standard prerequisites (editorial).}\begin{itemize}
\item Context span 3993-4008: P = k[u, u\textasciicircum{}\{-1\}, y, x\_n for n in Z] with the Z-action generated by tau(u) = u, tau(y) = u y, tau(x\_n) = x\_\{n+1\}; V(I\_d) with I\_d = ((1 - u\textasciicircum{}d) y, x\_n - x\_\{n+d\}) is the fix locus of tau\textasciicircum{}d; localizing at S generated by y and all 1 - u\textasciicircum{}d makes Z act freely on U = Spec(S\textasciicircum{}\{-1\}P); X = U/Z is the quotient algebraic space (Spaces, definition-quotient)
\item Context span 4010-4018: the primes p\_n = (x\_n, x\_\{n+1\}, ...) in S\textasciicircum{}\{-1\}P satisfy tau(p\_n) = p\_\{n+1\}, so they define points xi\_n of U with one common image x in X; the specializations xi\_n leads to xi\_\{n-1\} hold in U; the context concludes U -> X is an example of the promised type
\end{itemize}\medskip
\ifdefined\hypertarget\hypertarget{sourcepointer-1}{}\fi
\noindent\textit{Source reference.} Original definition 02Z3: sources/stacks-chow/tags.txt, source lines 3898--3898. Prerequisite pointer invoked at 4045; it names a standard fact of a sibling chapter that the counted argument uses, and every step of the counted proof built on it is retained verbatim in the delivered body.\par\medskip
To give an example, let $k$ be a field.
Set
$$
P = k[u, u^{-1}, y, \{x_n\}_{n \in \mathbf{Z}}].
$$
Consider the action
of $\mathbf{Z}$ on $P$ by $k$-algebra maps generated by the automorphism
$\tau$ given by the rules $\tau(u) = u$, $\tau(y) = uy$, and
$\tau(x_n) = x_{n + 1}$. For $d \geq 1$ set
$I_d = ((1 - u^d)y, x_n - x_{n + d}, n \in \mathbf{Z})$.
Then $V(I_d) \subset \Spec(P)$ is the fix point locus of $\tau^d$.
Let $S \subset P$ be the multiplicative subset generated by $y$ and
all $1 - u^d$, $d \in \mathbf{N}$. Then we
see that $\mathbf{Z}$ acts freely on $U = \Spec(S^{-1}P)$.
Let $X = U/\mathbf{Z}$ be the quotient algebraic space, see
Spaces, Definition \ref{spaces-definition-quotient}.

\medskip\noindent
Consider the prime ideals $\mathfrak p_n = (x_n, x_{n + 1}, \ldots)$ in
$S^{-1}P$. Note that $\tau(\mathfrak p_n) = \mathfrak p_{n + 1}$.
Hence each of these define point $\xi_n \in U$ whose image in $X$ is
the same point $x$ of $X$. Moreover we have the specializations
$$
\ldots \leadsto \xi_n \leadsto \xi_{n - 1} \leadsto \ldots
$$
We conclude that $U \to X$ is an example of the promised type.

\begin{lemma}
\label{lemma-specializations-fibre-etale}
There exists an \'etale morphism of algebraic spaces $f : X \to Y$
and a nontrivial specialization of points $x \leadsto x'$ in $|X|$ with
$f(x) = f(x')$ in $|Y|$.
\end{lemma}

\begin{proof}
See discussion above.
\end{proof}

\end{document}
